\documentclass[a4paper,12pt]{article}
\usepackage[utf8]{inputenc}
\usepackage[T2A]{fontenc}
\usepackage{cmap}
\usepackage[russian]{babel}
\usepackage{amsmath, amssymb, amsfonts}
\usepackage{graphicx}
\usepackage{float}   % для [H]
\usepackage{array}
\usepackage{geometry}
\geometry{top=2cm, bottom=2cm, left=2.5cm, right=2.5cm}

\title{Генератор истинно случайных чисел на основе джиттера в троичной логике}
\author{Р. Х. Латыпов, Е. Л. Столов\\
  Казанский федеральный университет, Казань, Россия}
\date{}

\begin{document}

\maketitle

\begin{abstract}
В данной работе представлено новое семейство генераторов, производящих истинно случайные равномерно распределённые числа в троичной логике. Генератор состоит из нескольких идентичных троичных логических комбинационных элементов, соединённых в кольцо. Предполагается, что каждый элемент имеет случайное время задержки, распределённое по экспоненциальному закону; все задержки независимы. Теория генератора основана на уравнениях Эрланга. Генератор может использоваться для тестирования в различных системах. Обсуждаются свойства многомерных случайных векторов, порождаемых генератором.
\end{abstract}

\section{Введение}

Генераторы случайных чисел (ГСЧ) важны в различных областях науки и техники, включая тестирование схем, стохастическое моделирование, генерацию криптографических ключей, случайную инициализацию переменных в криптографических протоколах и безопасность связи [1, 2]. Генератор истинно случайных чисел (TRNG) — это тип ГСЧ, который производит непредсказуемые данные, в отличие от псевдослучайного генератора, который даёт лишь статистически сильные (но полностью предсказуемые) последовательности. В последнее время было изобретено много типов TRNG. Большинство из них основаны на явлении, известном как «джиттер» в электронных схемах, — разновидности нестабильности сигналов, создаваемых электронным устройством [3]. Пример такого генератора, использующего нестабильность тактового генератора, управляющего D-триггером, представлен в [4]. Другая идея использует нестабильность выходов комбинационной схемы с обратными связями [5]. Простейший пример TRNG, основанный на этом изобретении, показан на рис. 1. Эта идея использовалась многими авторами в двоичном случае. Некоторые ссылки на оригинальные работы можно найти в [5–7]. При определённых предположениях о процессе можно разработать теорию джиттера и создать схему, работающую как генератор случайных значений. Недавно было предложено использовать генераторы с джиттером, построенные на основе перепрограммируемых цифровых схем или конструкций, основанных на метастабильности [6, 8]. В обоих подходах джиттер используется для генерации битовых последовательностей [5, 6]. Несмотря на существование ряда аппаратных TRNG, задача поиска эффективного метода генерации истинно случайных чисел, реализуемого только с помощью логических вентилей, остаётся актуальной.

Основными недостатками двоичных элементов являются проблемы соединений и выводов. С другой стороны, в последние годы всё больше изучаются троичные схемы и их возможная физическая реализация. Подробное обсуждение можно найти в [9–11]. Более того, были разработаны наноэлектронные устройства нового поколения (такие как резонансные туннельные приборы и квантово-точечные клеточные автоматы) [10, 11]. Такие схемы часто являются многоуровневыми и могут оказаться более совершенными в многоразрядных цепях. Разработка троичных логических устройств представляется перспективной для повышения скорости арифметических операций; следовательно, разработка таких схем становится актуальной задачей. В нашей статье мы представляем новое семейство TRNG, которое может быть реализовано с использованием только троичных логических вентилей. Основная цель работы — теоретическое исследование поведения таких устройств.

\begin{figure}[H]
  \centering
  \framebox[4cm]{\parbox{4cm}{\centering Рис. 1. Инвертор в режиме джиттера.}}
  \caption{Инвертор в режиме джиттера}
\end{figure}

Оставшаяся часть статьи организована следующим образом. В разделе II приведены базовая схема генератора и её математическая теория; в разделе III исследуются статистические свойства генератора; раздел IV посвящён реализации генератора; выводы даны в разделе V.

На протяжении всей статьи при работе с матрицами используются следующие обозначения:

(i) \(A[i,j]\) — элемент матрицы \(A\), стоящий в \(i\)-й строке и \(j\)-м столбце.  
(ii) \(A[i,*]\) — \(i\)-я строка матрицы \(A\), \(A[*,j]\) — \(j\)-й столбец этой матрицы.  
(iii) \(\theta\) — нулевая матрица.

\section{Предварительные сведения и базовая схема}

В данной работе предлагается и анализируется TRNG, основанный на сбалансированной троичной логике. Сбалансированная троичная логика — это частный случай троичной логики, где цифры принимают значения \(-1, 0, 1\). Мы будем работать с двухвходовыми одновыходными функциями или вентилями. Это означает, что существует \(3^2 = 9\) возможных входов и \(3^9 = 19383\) возможных вентилей. Рассмотрим схему, показанную на рис. 2. Здесь \(F\) обозначает троичный комбинационный вентиль, \(c = F(a,b)\), \(a,b,c \in \{-1,0,1\}\). В общем случае схема состоит из \(N\) вентилей. Перенумеруем все элементы числами из интервала \([0, N-1]\). Выход элемента с номером \(k\) подключён к одному из входов этого же элемента и ко входу элемента с номером \(k-1 \bmod N\) (см. рис. 2). После того как хотя бы один входной сигнал изменяется, выходной сигнал также может измениться, но это требует времени задержки \(DT\). В статье приняты следующие предположения:

(i) \(DT\) — стохастическая величина, имеющая экспоненциальное распределение для всех элементов с одинаковым параметром \(T\). Это означает, что \(P(DT < d) = 1 - \exp(-Td)\).  
(ii) В любой момент времени только один элемент может изменить свой выход; все эти события независимы.

В дальнейшем будем считать \(T = 1\), это не ограничивает общности рассуждений. Все ограничения, налагаемые на систему, позволяют применить теорию Эрланга [12] для описания функционирования генератора. Пусть \(s_k(t)\) — выходной сигнал элемента с номером \(k\) в момент времени \(t\). Обозначим через \(\mathbf{S}(t) = \langle s_0(t), \ldots, s_{N-1}(t) \rangle\) состояние TRNG в момент времени \(t\). Таким образом, система имеет \(M = 3^N\) состояний, которые также индексируются числами от 0 до \(M-1\). Для каждого \(n \in [0, M-1]\) пусть \(i_1, \ldots, i_m\) — список номеров состояний, из которых возможен переход в состояние \(\mathbf{S}_n\) в результате изменения выходного сигнала одного из элементов. Этот список зависит от \(n\). Составим систему дифференциальных уравнений типа Эрланга, описывающих динамику TRNG [12]. Пусть \(P_n(t)\) — вероятность того, что TRNG находится в состоянии с номером \(n\) в момент времени \(t\). Вероятность \(P_n(t + \Delta t)\) равна сумме вероятностей: \((1 - \Delta t)^N P_n(t)\) — ни один элемент не изменил выход, так что TRNG не изменил состояния; и \(\Delta t \sum_{i,k} P_{i,k}\) — ровно один элемент изменил выход и TRNG перешёл в состояние \(\mathbf{S}_n\). Здесь использованы очевидные приближения: \(\exp(-\Delta t) \sim 1 - \Delta t\) и \(1 - \exp(-\Delta t) \sim \Delta t\) для малых \(\Delta t\). Устремляя \(\Delta t\) к нулю, получаем для каждого \(n \in [0, M-1]\)

\begin{equation}
\frac{dP_n}{dt} = -N P_n(t) + \sum_{k=1}^{m} P_{i_k}(t), \quad (1)
\end{equation}

где \(m, i_1, \ldots, i_m\) зависят от \(n\). Это уравнение типа Эрланга, обычно используемое в теории массового обслуживания.

\subsection{Выбор функции \(F\)}

Уравнение (1) справедливо для любой функции \(F\), упомянутой выше. Теперь необходимо наложить некоторые ограничения на эту функцию, чтобы гарантировать определённые свойства TRNG. Прежде всего, TRNG не должен иметь устойчивых состояний. Это означает, что для любого \(n\) выполняется \(\forall t > 0 \; P_n(t) < 1\). Предположим, что выполняется следующее условие:

\begin{equation}
c = F(a,b), \quad \forall (a,b) \quad c \neq a, b. \quad (2)
\end{equation}

Из (2) следует, что любой элемент изменит свой выход после времени задержки, так что TRNG не имеет устойчивых состояний. Найдём все функции, обладающие свойством (2). Прежде всего, если \(a \neq b\), то \(F(a,b)\) определяется однозначно, и поэтому \(F(a,b) = F(b,a)\). Последнее означает, что нам нужно определить только \(F(-1,-1)\), \(F(0,0)\) и \(F(1,1)\). Пусть \(\sigma\) — произвольная перестановка значений \(-1,0,1\). Функции \(F(a,b)\) и \(\sigma(F(\sigma(a), \sigma(b)))\) можно считать одинаковыми. Две функции, удовлетворяющие (2), приведены в таблице 1.

\begin{table}[H]
\centering
\begin{tabular}{|c|c|c|c|}
\hline
Имя & \(F(0,0)\) & \(F(1,1)\) & \(F(-1,-1)\) \\
\hline
\(F_1\) & 1 & -1 & 0 \\
\(F_2\) & 1 & -1 & 1 \\
\hline
\end{tabular}
\caption{Функции, пригодные для построения генератора}
\label{tab:functions}
\end{table}

\paragraph{Предложение 1.} Существует только две различные троичные функции, удовлетворяющие (2).

Доказательство предложения приведено в приложении. Далее будем считать, что \(F = F_1\).

\subsection{Динамика состояний генератора}

После того как тип \(F\) определён, можно найти все параметры в (1). Используя матричную форму, перепишем систему следующим образом:

\begin{equation}
\frac{d\mathbf{P}}{dt} = \mathrm{Matr} \cdot \mathbf{P}, \quad (3)
\end{equation}

где \(\mathbf{P} = \langle P_0, P_1, \ldots, P_{M-1} \rangle^T\). Поскольку размер матрицы \(\mathrm{Matr}\) быстро растёт с увеличением числа элементов в генераторе, мы ограничимся случаями \(N=1\), \(N=2\) или \(N=3\) в примерах. Прежде всего, необходимо проиндексировать все состояния TRNG. Пусть \(\mathbf{S} = \langle s_0, s_1, \ldots, s_{N-1} \rangle\) — состояние TRNG, \(s_k \in \{-1,0,1\}\). Индекс состояния \(\mathbf{S}\) — это число

\begin{equation}
\operatorname{Ind}(\mathbf{S}) = \sum_{k=0}^{N-1} 3^k (s_k + 1). \quad (4)
\end{equation}

Например, для \(N=2\) матрица \(\mathrm{Matr}\) представлена в (5):

\[
\mathrm{Matr} = \left( \begin{array}{cccccccccc}
-2 & 0 & 0 & 0 & 0 & 0 & 0 & 0 & \dots & 0 \\
1 & -2 & 1 & 0 & 0 & 0 & 0 & 1 & 0 & 0 \\
0 & 1 & -2 & 0 & 0 & 1 & 0 & 0 & 1 & 0 \\
1 & 0 & 0 & -2 & 0 & 1 & 1 & 0 & 0 & 0 \\
0 & 0 & 0 & 0 & -2 & 0 & 0 & 0 & 0 & 0 \\
0 & 0 & 1 & 1 & 1 & -2 & 0 & 0 & 0 & 0 \\
0 & 0 & -2 & 0 & 0 & -2 & 1 & 1 & 0 & 0 \\
0 & 1 & 0 & 0 & 1 & 0 & 1 & -2 & 0 & 0 \\
0 & 0 & 0 & 0 & 0 & 0 & 0 & 1 & -2 & 0
\end{array} \right) \quad (5)
\]

Поясним смысл элементов этой матрицы. Согласно (4), состояние \(\langle 0,1 \rangle\) имеет индекс 7. В строке с номером 7 (начиная с нуля) матрица \(\mathrm{Matr}\) содержит 1 в позициях 1, 4, 6. Из (4) следует, что эти индексы соответствуют состояниям \(\langle 0,-1 \rangle\), \(\langle 0,0 \rangle\) и \(\langle -1,1 \rangle\). Это означает, что TRNG может перейти из этих состояний в состояние 7. Действительно, предположим, что TRNG находится в состоянии 6. Элемент 0 имеет на своих входах \(-1\) от элемента 0 и \(1\) от элемента 1. В соответствии с определением \(F_1\), он меняет свой выход на 0, переходя в состояние 7. Все остальные случаи анализируются аналогично.

Формально, имея матрицу \(\mathrm{Matr}\) и дифференциальные уравнения (3), можно найти вектор \(\mathbf{P}(t)\) для любого \(t\). Наша цель — получить свойства этого вектора без поиска точного решения системы.

\section{Асимптотические свойства \(\mathbf{P}(t)\)}

Строки с индексами 0, 4 и 8 в (5) не содержат единиц. Это означает, что переходов TRNG в состояния с такими индексами нет; иными словами, эти состояния недостижимы.

\paragraph{Предложение 2.} Состояния вида \(\langle x, x, \ldots, x \rangle\), где \(x \in \{-1,0,1\}\), являются единственными недостижимыми состояниями генератора при \(N > 1\).

Доказательство предложения приведено в приложении.

В общем случае исключим из рассмотрения состояния вида \(\langle x, x, \ldots, x \rangle\). Это означает, что нужно исключить три элемента из вектора \(\mathbf{P}\) и три строки и три столбца из матрицы \(\mathrm{Matr}\) в (3). Сохраним прежние обозначения для новых \(\mathbf{P}\) и \(\mathrm{Matr}\). Поскольку матрица \(\mathrm{Matr}\) в (3) постоянна, решение уравнения можно представить в виде

\begin{equation}
\mathbf{P}(t) = \exp(\mathrm{Matr} \cdot t) \cdot \mathbf{P}(0), \quad (6)
\end{equation}

где \(\mathbf{P}(0)\) — стохастический вектор, определяющий начальное состояние TRNG [14]. Согласно определению,

\begin{equation}
\mathrm{Matr} = Q - N \cdot I, \quad (7)
\end{equation}

где \(I\) — единичная матрица, а \(Q\) — бинарная матрица (содержит 1 или 0). Обе матрицы имеют размер \(M' \times M'\), \(M' = M - 3\).

\paragraph{Предложение 3.} Пусть \(e_1, e_2, \ldots, e_{M'}\) — все собственные значения \(Q\) и
\[
|e_1| \ge |e_2| \ge \cdots \ge |e_{M'}|. \quad (8)
\]
Тогда \(e_1 = N\).

Доказательство предложения приведено в приложении.

Отсюда следует, что \(m_1 = 0\) является собственным значением для \(\mathrm{Matr}\), а для всех остальных собственных значений \(m_2, \ldots, m_{M'}\) этой матрицы выполняется неравенство \(\operatorname{Re}(m_j) \le 0\), \(j = 2, \ldots, M'\). Предположим, что имеет место строгое неравенство
\[
\operatorname{Re}(m_j) < 0, \quad j = 2, \ldots, M'. \quad (9)
\]
Это означает, что матрица \(E(t) = \exp(t \cdot \mathrm{Matr})\) имеет единственный корень 1, а все остальные корни этой матрицы имеют модуль меньше 1. При \(t \to \infty\) матрица \(E(t)\) сходится к устойчивой матрице. То же самое справедливо и для вектора \(\mathbf{P}(t)\), и
\[
\frac{d\mathbf{P}(t)}{dt} \to 0, \quad t \to \infty.
\]
Следуя [12], можно найти стохастический вектор \(\bar{\mathbf{P}} = \mathbf{P}(\infty)\), решая систему
\begin{equation}
Q \cdot \bar{\mathbf{P}} = N \bar{\mathbf{P}}. \quad (10)
\end{equation}

\subsection{Случай \(N=1\)}

Существует три состояния TRNG: \(\langle -1 \rangle\), \(\langle 0 \rangle\), \(\langle 1 \rangle\) и
\[
\mathrm{Matr} = \begin{pmatrix}
-1 & 0 & 1 \\
1 & -1 & 0 \\
0 & 1 & -1
\end{pmatrix}.
\]
Собственные значения \(\mathrm{Matr}\) равны \(0, -1.5, -1.5\), таким образом, (9) выполняется, и \(\bar{\mathbf{P}} = \langle 1/3, 1/3, 1/3 \rangle\). Это означает, что устойчивый вектор \(\bar{\mathbf{P}}\) имеет равномерное распределение.

\subsection{Случай \(N=2\)}

Здесь собственные значения \(\mathrm{Matr}\) равны \(0, -1, -4, -3, -3, -1\), а устойчивый вектор равен \(\langle 0.17, 0.17, 0.17, 0.17, 0.17, 0.17 \rangle\); следовательно, и в этом случае имеем равномерное распределение.

\subsection{Случай \(N=3\)}

В этом случае (9) выполняется, но устойчивый вектор \(\bar{\mathbf{P}}\) имеет компоненты (в компактной записи)
\[
\begin{array}{llllll}
0.037 & 0.037 & 0.037 & 0.037 & 0.074 & 0.037 \\
0.037 & 0.037 & 0.037 & 0.037 & 0.037 & 0.037 \\
0.037 & 0.074 & 0.037 & 0.037 & 0.037 & 0.074 \\
0.037 & 0.037 & 0.037 & 0.037 & 0.037 & 0.037
\end{array}
\]
(всего 24 компоненты). Следовательно, мы имеем неравномерное финальное распределение состояний. Более того, существуют особые состояния \(\langle -1,0,1 \rangle\), \(\langle 1,-1,0 \rangle\), \(\langle 0,1,-1 \rangle\). Фактически это одно и то же состояние, поскольку каждое является результатом циклического преобразования других (см. рис. 2). Соответствующие строки в таблице содержат шесть единиц, тогда как другие строки содержат по три единицы.

\subsection{Общий случай}

Напомним [13], что матрица \(Q\) называется разложимой, если существует матрица перестановок \(Pr\) такая, что
\[
Pr^T \cdot Q \cdot Pr = \begin{pmatrix} A & 0 \\ B & C \end{pmatrix},
\]
где \(A, C\) — квадратные матрицы. Можно доказать, что в общем случае матрица \(Q\) в (7) является неразложимой. Из-за ограничения на объём сообщения мы опустим доказательство этого факта. Оно будет представлено позже в статье, посвящённой этой проблеме. Известно ([13]), что для неразложимых матриц выполняется (9). Это означает, что метод расчёта распределения состояний, представленный выше, применим для любого числа элементов в TRNG. Единственная проблема — найти собственные векторы матрицы большого размера. Примеры показывают, что устойчивые распределения \(\bar{\mathbf{P}}\) не являются равномерными при \(N > 2\). Поскольку соответствующие устойчивые векторы имеют большой размер, мы представляем компоненты вектора в графическом виде (см. рис. 3). Это означает, что существует некоторая зависимость между различными выходами TRNG. С другой стороны, из-за симметрии схемы выходной сигнал любого элемента имеет равномерное распределение.

\begin{figure}[H]
  \centering
  \framebox[6cm]{\parbox{6cm}{\centering Рис. 3. Компоненты устойчивого вектора для \(N=4\) (график).}}
  \caption{Компоненты устойчивого вектора}
\end{figure}

\section{Реализация троичного генератора}

Стандартная реализация генератора представлена на рис. 4. Линия тактовых импульсов обозначает серию импульсов с интервалом \(D\). Необходимо выбрать \(D\) достаточно большим, чтобы время было достаточным для перехода TRNG в устойчивое состояние. В реальной ситуации невозможно получить устойчивое распределение за конечное время; поэтому можно достичь лишь приближения к устойчивому распределению.

\begin{figure}[H]
  \centering
  \framebox[4cm]{\parbox{4cm}{\centering Рис. 4. Реализация. Здесь gen — TRNG, выходы которого управляются тактовым сигналом.}}
  \caption{Реализация генератора}
\end{figure}

Небольшая нестабильность \(D\) не влияет на это распределение. Для этого необходимо оценить разницу между устойчивой матрицей \(\mathrm{Stab}\) размера \(M' \times M'\) вида
\[
\mathrm{Stab} = (\bar{\mathbf{P}}, \bar{\mathbf{P}}, \ldots, \bar{\mathbf{P}})
\]
и матрицей \(\mathrm{Astab}(D) = \exp(D \cdot \mathrm{Matr})\). Пусть \(\mathbf{E} = \langle 1,1,1,\ldots,1 \rangle\). Согласно определению, \(\mathbf{E} \cdot \mathrm{Matr} = \Theta = \langle 0,0,\ldots,0 \rangle\); следовательно, \(\mathbf{E} \cdot \mathrm{Matr}^k = \Theta\) для любого натурального \(k\), и \(\mathbf{E} \cdot \mathrm{Astab}(D) = \mathbf{E}\). Другими словами, сумма элементов в любом столбце \(\mathrm{Astab}(D)\) равна 1. Пусть
\[
\delta = \max_{k} \left| \bar{\mathbf{P}}[k] - \mathrm{Astab}(D)[*,k] \right|, \quad (12)
\]
где для вектора \(\mathbf{q} = \langle a_1, \ldots, a_n \rangle\) \(|\mathbf{q}| = \max_k |a_k|\). Выберем \(\delta\) как расстояние между двумя матрицами. Некоторые результаты расчётов представлены в таблице 2.

\begin{table}[H]
\centering
\begin{tabular}{|c|c|c|c|c|}
\hline
\(N\) \textbackslash \(D\) & 2 & 4 & 8 & 16 \\
\hline
1 & \(3.1\cdot 10^{-2}\) & \(1.6\cdot 10^{-3}\) & \(3.9\cdot 10^{-6}\) & \(2.1\cdot 10^{-7}\) \\
2 & \(4.6\cdot 10^{-2}\) & \(6.1\cdot 10^{-3}\) & \(1.1\cdot 10^{-4}\) & \(3.8\cdot 10^{-8}\) \\
3 & \(3.6\cdot 10^{-2}\) & \(6.9\cdot 10^{-3}\) & \(2.4\cdot 10^{-4}\) & \(3.2\cdot 10^{-7}\) \\
4 & \(2.1\cdot 10^{-2}\) & \(3.7\cdot 10^{-3}\) & \(1.4\cdot 10^{-4}\) & \(2.6\cdot 10^{-7}\) \\
5 & \(9.8\cdot 10^{-3}\) & \(2.1\cdot 10^{-3}\) & \(9.2\cdot 10^{-5}\) & \(2.0\cdot 10^{-7}\) \\
6 & \(6.0\cdot 10^{-3}\) & \(9.4\cdot 10^{-4}\) & \(3.6\cdot 10^{-5}\) & \(7.1\cdot 10^{-8}\) \\
\hline
\end{tabular}
\caption{Расстояние между устойчивой матрицей и \(\mathrm{Astab}(D)\) в зависимости от \(D\)}
\end{table}

Как уже упоминалось, недостатком предложенного TRNG является отсутствие трёх специальных сигналов вида \(\langle x,x,\ldots,x \rangle\), \(x \in \{-1,0,1\}\), которые не выдаются на выходах схемы. Можно предложить два метода преодоления этого недостатка. Первый — использование набора независимых TRNG (рис. 5). Эта схема выдаёт равномерно распределённые сигналы.

\begin{figure}[H]
  \centering
  \framebox[4cm]{\parbox{4cm}{\centering Рис. 5. Генератор на основе независимых элементов.}}
  \caption{Генератор на основе независимых элементов}
\end{figure}

Каждый TRNG на рис. 5 содержит один элемент, но можно использовать TRNG с произвольным числом элементов, хотя в таком случае будет использоваться выход только одного элемента. Существует другое решение проблемы отсутствия трёх специальных состояний на выходах TRNG. Предположим, у нас есть схема, представленная на рис. 6. Выход элемента 0 опущен. Из-за циклической структуры TRNG можно опустить любой из выходов, но распределение выходных сигналов будет одинаковым. В результате на выходах TRNG может возникнуть любой троичный вектор, и распределение этих векторов можно найти на основе предыдущих расчётов. Например, для \(N=3\), если используются только два выхода, получаем вероятности (см. (11)):  
(i) состояния \(\langle x,x \rangle\) — 0.074;  
(ii) состояния \(\langle 0,1 \rangle\), \(\langle -1,0 \rangle\), \(\langle 1,-1 \rangle\) — 0.148;  
(iii) все остальные состояния — 0.111.  
Таким образом, распределение состояний немного ближе к равномерному, чем в исходном примере.

\begin{figure}[H]
  \centering
  \framebox[4cm]{\parbox{4cm}{\centering Рис. 6. Генератор с опущенным выходом.}}
  \caption{Генератор с опущенным выходом}
\end{figure}

\section{Заключение}

Предположение о том, что времена задержки в различных элементах являются независимыми событиями, является естественным, тогда как ограничение на вид распределения значительно сильнее. Из-за симметрии схемы, если время задержки для всех элементов имеет одинаковое распределение, то выход каждого элемента будет равномерным независимо от распределения. Тем не менее существование устойчивого распределения и его вид, если оно существует, требуют дополнительных исследований.

\appendix
\section{Доказательство предложения 1}

Пусть \(F(a,b)\) удовлетворяет (2) и \(F(0,0)=c_1\), \(F(1,1)=c_2\), \(F(-1,-1)=c_3\). Имеются следующие альтернативы:

(i) \(c_1, c_2, c_3\) — перестановка \(0,1,-1\);  
(ii) ровно два элемента из множества \(c_1, c_2, c_3\) равны.

Рассмотрим первый случай. Пусть \(c_1 = 1\). Так как все \(c_i\) различны и \(c_3 \neq -1\), имеем \(c_3 = 0\). Тогда \(F = F_1\). Если \(c_1 = -1\), то \(c_2 = 0\), \(c_3 = 1\). Используя перестановку \(\sigma(0)=0, \sigma(1)=-1, \sigma(-1)=1\), преобразуем \(F\) к \(F_1\).

Теперь рассмотрим второй случай. Рассмотрим различные возможные распределения значений \(c_1, c_2, c_3\).

(i) \(c_1 = 1, c_2 = 0, c_3 = 1\). Перестановка \(\sigma(0)=-1, \sigma(1)=1, \sigma(-1)=0\) преобразует \(F\) в \(F_2\).  
(ii) \(c_1 = 1, c_2 = 0, c_3 = 0\). Перестановка \(\sigma(0)=1, \sigma(1)=-1, \sigma(-1)=0\) преобразует \(F\) в \(F_2\).

Если \(c_1 = -1\), то перестановка \(\sigma(0)=0, \sigma(1)=-1, \sigma(-1)=1\) сводит дело к предыдущей ситуации.

\section{Доказательство предложения 2}

Сначала покажем, что \(\mathbf{S} = \langle x,x,\ldots,x \rangle\) является недостижимым состоянием. Поскольку за один переход может измениться только один выход, любое состояние-кандидат для перехода в \(\mathbf{S}\) должно иметь вид \(\langle x,x,\ldots,x,y,\ldots,x \rangle\). Любой элемент в схеме имеет на входах либо \(\langle x,x \rangle\), либо \(\langle x,y \rangle\), либо \(\langle y,x \rangle\). Во всех этих случаях элемент может изменить свой выход на \(x\) в соответствии с определением \(F\). С другой стороны, если \(\mathbf{S} = \langle \ldots,x,y,\ldots \rangle\), где \(x \neq y\), то состояние \(\langle \ldots,x,z,\ldots \rangle\), где \(z \neq x\) и \(z \neq y\), может перейти в \(\mathbf{S}\).

\section{Доказательство предложения 3}

Согласно определению, элемент \(Q[i,j]\) в строке с номером \(i\) и столбце с номером \(j\) равен 1 тогда и только тогда, когда генератор может перейти из состояния с номером \(j\) в состояние с номером \(i\). Поскольку генератор имеет \(N\) элементов, каждый из них может изменить свой выход, и все новые состояния различны. Отсюда следует, что для каждого \(j\)
\[
\sum_i Q[i,j] = N. \quad (13)
\]
Это означает, что \(Q^T / N\) является стохастической матрицей, следовательно, \(e_1 = N\), и неравенство (8) выполняется.

\begin{thebibliography}{99}
\bibitem{1} Asmussen S and Glynn P W 2007 \textit{Stochastic Simulation: Algorithms and Analysis} (New York: Springer Verlag) p 476
\bibitem{2} Ferguson N, Schneie B, and Kohno T 2010 \textit{Cryptography Engineering: Design Principles and Practical Applications} (Indianapolis: Wiley) p 384
\bibitem{3} Horowitz P and Hill W 1980 \textit{The art of electronics} (Cambridge: Cambridge University Press) p 1125
\bibitem{4} Petrie C S and Connelly J A 1996 A noise-based IC random number generator for applications in cryptography \textit{Proc. IEEE Int. Symp. Circuits and Systems} (Atlanta) vol 4 (New York: IEEE Press) pp 324-327
\bibitem{5} Golic J D 2006 New Methods for Digital Generation and Postprocessing of Random Data \textit{IEEE Trans. Comput.} \textbf{55} 1217-29
\bibitem{6} Sunar B, Martin W J and Stinson D R 2007 A provably secure true random number generator with built-in tolerance to active attacks \textit{IEEE Trans. Comput.} \textbf{56} 109-19
\bibitem{7} Kuznetsov V, Pesoshin V and Stolov E 2008 Markov model of a digital stochastic generator \textit{Automation and Remote Control} \textbf{69} 1504-09
\bibitem{8} Wieczorek P and Golofit K 2014 Dual-Metastability Time-Competitive True Random Number Generator \textit{IEEE Trans. Circuits and Systems} \textbf{61} 134-45
\bibitem{9} Wu X W and Prosser F P 1990 CMOS ternary logic circuits \textit{IEE Proc. Circuits, Devices and Systems} \textbf{137} 21-27
\bibitem{10} Gaikwad V N and Deshmukh P R 2015 Design of CMOS ternary logic family based on single supply voltage \textit{Proc. IEEE Int. Conf. on Pervasive Computing} (Sydney) vol 9 (New York: IEEE Press) pp 1-6
\bibitem{11} Lisa N J and Babu H H 2015 Design of a Compact Ternary Parallel Adder/Subtractor Circuit in Quantum Computing \textit{Proc. IEEE Int. Symp. on Circuits and Systems} (Lisbon) vol 28 (New York: IEEE Press) pp 2145-2148
\bibitem{12} Kleinrock L 1975 \textit{Queueing Systems: Volume I Theory} (New York: Wiley-Interscience) p 417
\bibitem{13} Marcus M and Mink H 1964 \textit{A survey of matrix theory and matrix inequalities} (Boston: Allys and Bacon) p 232
\bibitem{14} Bellman R 1960 \textit{Introduction to matrix analysis} (New York: Macgrow-Hill) p 365
\end{thebibliography}

\end{document}